\section{Communication-Graph Representation and E-GAT-IDS}
\label{sec:method}

\subsection{Graph construction}

We represent the traffic as a directed multigraph
$G = (V, E, X, Z)$ built as follows.

\emph{Nodes.} $V = \{v_1,\dots,v_N\}$ contains one node per unique IP endpoint
observed as either a source or a destination, so $N = 24$ for both flow
abstractions. The IP string is used for identity only and never enters a
feature vector.

\emph{Edges.} Each flow record $f$ contributes one directed edge
$e_f = (v_{s(f)}, v_{d(f)})$. Flows sharing an ordered endpoint pair are
retained as \emph{parallel edges} rather than collapsed, because the
classification target belongs to the individual record: collapsing would
destroy the task. The uni-flow graph therefore has 495{,}525 edges over 106
distinct ordered pairs, a mean parallel multiplicity of roughly 4{,}675.
Figure~\ref{fig:flow2graph} illustrates the construction. A line-graph
alternative, in which flows become nodes and shared endpoints become edges, is
implemented for completeness but is not used for the headline results; the
exact line graph of a multigraph with this multiplicity is quadratic in the
multiplicity and does not fit in memory.

\emph{Edge attributes.} $Z \in \mathbb{R}^{|E| \times D}$ holds the
preprocessed flow attributes, $D = 27$ for uni-flow and $D = 40$ for bi-flow
after the deterministic derivations of Section~\ref{sec:methodology}.

\emph{Node attributes.} $X \in \mathbb{R}^{N \times D_n}$ holds
behaviour-derived features, $D_n = 25$: in- and out-degree, distinct peer
counts in each direction, flow and packet and byte volumes in each direction,
mean packet length and mean and standard deviation of inter-arrival time in
each direction, protocol diversity, source- and destination-port diversity,
destination-port entropy, and three ratio features. We use these rather than a
learned embedding indexed by endpoint deliberately. An embedding table keyed on
IP address is free to memorise that one particular address is the attacker,
which on a fixed-address testbed is a shortcut that inflates results without
conferring capability. Behaviour aggregates at least require the model to
express \emph{what} an endpoint does.

Node attributes are aggregated over the \emph{training} flows only. A test
flow's attributes therefore never influence any node's representation, even
indirectly. We report the transductive alternative---aggregating over all
flows---as an explicitly labelled contrast, and treat any gain it produces as
leakage rather than capability.

\subsection{Message passing}

Let $h_i^{(0)} = \sigma(W_0 X_i)$ be the projected node state. A layer updates
node states by attention-weighted aggregation over incoming edges,
\begin{equation}
h_i^{(l+1)} = \sigma\!\Big( \sum_{j \in \mathcal{N}(i)}
   \alpha_{ji}\, m_{ji}^{(l)} \Big),
\label{eq:mp}
\end{equation}
where $\mathcal{N}(i)$ are the endpoints with an edge into $v_i$,
$\alpha_{ji}$ is the attention weight and $m_{ji}^{(l)}$ the message. The
adjacency used in Eq.~\eqref{eq:mp} contains training edges only, so a test
flow is scored against node states never informed by test traffic.

Because a handful of hub endpoints---the broker and the attacker---carry
hundreds of thousands of incident edges, we sample at most $K$ incoming
training edges per destination node, redrawn each optimisation step and fixed
under a recorded seed at evaluation. This is neighbour sampling in the
GraphSAGE sense~\cite{hamilton2017graphsage} and not a change to the model
definition; $K = 2000$ during training and $K = 4000$ at evaluation.

\subsection{Edge-aware attention}

The standard graph attention logit~\cite{velickovic2018gat} is computed from
the two endpoint states alone,
\begin{equation}
e_{ij} = \mathrm{LeakyReLU}\big(\mathbf{a}^{\top}
  [\,W h_i \,\Vert\, W h_j\,]\big).
\label{eq:gat}
\end{equation}
Under Eq.~\eqref{eq:gat} every parallel edge between the same ordered pair
receives the same attention weight, regardless of how the flows differ. Since
the flows between one sensor and the broker include both benign publishes and,
in the brute-force capture, attack traffic, this is a representational
limitation and not merely an inefficiency.

The edge-aware variant admits a projection of the flow attributes into the
logit,
\begin{equation}
e_{ij} = \mathrm{LeakyReLU}\big(\mathbf{a}^{\top}
  [\,W h_i \,\Vert\, W h_j \,\Vert\, U z_{ij}\,]\big),
\qquad
\alpha_{ij} = \operatorname*{softmax}_{j}\big(e_{ij}\big),
\label{eq:egat}
\end{equation}
so that a flow's own characteristics modulate how strongly it is attended to.
Figure~\ref{fig:attention} contrasts the two. This placement is that of
EGAT~\cite{chen2021egat}; we claim no novelty for it and treat it as the
mechanism under test. The message may likewise carry the projected attributes,
$m_{ij} = W h_i + U z_{ij}$, which is separately switchable.

Attention uses $H$ heads with per-head dimension $d_h$, concatenated across
heads, with layer normalisation, ELU activation, dropout and a residual
connection around each layer.

\subsection{Edge classification}

After $L$ layers, each flow is classified from the concatenation of its two
endpoint states and its own attributes,
\begin{equation}
r_{ij} = [\, h_i \,\Vert\, h_j \,\Vert\, z_{ij} \,], \qquad
\hat{y}_{ij} = \operatorname{softmax}\big(\mathrm{MLP}(r_{ij})\big).
\label{eq:clf}
\end{equation}
Figure~\ref{fig:architecture} shows the pipeline. Retaining $z_{ij}$ in
Eq.~\eqref{eq:clf} means the graph models are never denied information the
tabular baselines have; the comparison then isolates what the graph
\emph{adds} rather than what withholding features costs. A variant with
$z_{ij}$ removed from Eq.~\eqref{eq:clf} is included in the ablation.

\subsection{The baseline family as a single code path}

Every graph baseline is this same architecture under different settings of
three switches---edge attributes in the attention logit, in the message, and in
the classifier---and a choice of aggregator. GCN uses symmetric-normalised
sum aggregation~\cite{kipf2017gcn} and GraphSAGE mean
aggregation~\cite{hamilton2017graphsage}, both topology-only; E-GraphSAGE is
mean aggregation with attributes in the message, which is the published
configuration~\cite{lo2022egraphsage}; GAT is attention without attributes;
EdgeGAT is attention with attributes in the message only; and E-GAT-IDS adds
them to the attention logit. Because the scaffolding, the training loop, the
metric path and the result writer are shared, a difference between two rows is
attributable to the switch that differs and not to an unrelated
implementation choice.

\subsection{Objective}

We use class-weighted cross entropy with inverse-frequency weights computed
from the training labels, and evaluate focal loss~\cite{lin2017focal} and
weight-free variants in the ablation. We do not apply
SMOTE~\cite{chawla2002smote} to the graph models: interpolating between flow
records would fabricate communication events that were never observed, which on
a graph whose edges are supposed to be observed traffic is not defensible. Where
resampling is used, existing edges are resampled so that every training example
remains a real flow. SMOTE \emph{is} applied inside the GA-EBT reproduction,
because the original specifies it.
